We close this section with a run of the function $\SATALG$.
This run illustrates the inner workings of the algorithm it implements.
Moreover, it allows us to recognize how it uses only polynomial space.

\begin{exa}
Let us suppose that we start $\SATALG$ in~\Cref{alg:start} with the node expression
\[
    \varphi = \tup{\child(1\land p)? =_\cmp \astep[b]} \land \tup{\astep[c]@_1(\neg p)? =_\cmp \astep[b]}  \land \neg\tup{\astep[b]\neq_\cmp @_1}.
\]   
The call to $\SATALG(\varphi)$ invokes $\SATALG(\Phi,\NOM)$ with $\Phi=[2{:}\varphi]$ and  $\NOM=3$.
We proceed to list the different steps in~\Cref{alg:sat} on this input.
Sometimes we group several steps together and omit some node expressions to make the presentation more succinct. 

\begin{description}
  \item[Initialization] 
    We set $\BRANCH=\Phi$ and $\STACK=[~]$.
    Notice that the sole root nominal in $\varphi$ is $1$.

  \item[Step 1] 
    We reach the `expand' block and apply rules of type 1 until no such rule is left untreated.
    The list $\BRANCH$ is updated during this process (the rules applied to obtain the node expressions are indicated on the right):
    \[
      \begin{array}{r@{\,}@{}lr}
        \BRANCH =
          & [
              2{:}\tup{\child(1\land p)? =_\cmp \astep[b]},
              2{:}\tup{\astep[c]@_1(\neg p)? =_\cmp \astep[b]},
              2{:}\neg\tup{\astep[b]\neq_\cmp @_1},
          ] + \dots +
              & (\land) \\
          & [
              \tup{@_2\child(1\land p)? =_\cmp @_2\astep[b]},
              \tup{@_2\astep[c]@_1(\neg p)? =_\cmp @_2\astep[b]},
              \neg\tup{@_2\astep[b]\neq_\cmp @_2@_1}
            ]
              & (\text{int}_1), (\text{int}_2)
      \end{array}
    \]

  \item[Step 2]
    We reach the `explore' block and choose to apply the rule of type 3 with premiss $\tup{@_2\child(1\land p)? =_\cmp @_2\astep[b]}$.
    Before the recursive call we get
    \[
      \Phi =
        [
          \neg\tup{@_2\astep[b]\neq_\cmp @_2@_1}, 2{:}\tup{\astep[a]}3, \tup{@_3(1\land p)? =_\cmp @_2\astep[b]}
        ]
    \]

  \item[Step 3] 
    We call $\SATALG(\Phi,\NOM)$ recursively for a first time, with $\NOM=4$. 
    In the `expand' block of the first recursive call, we apply all rules of type 1 and get
    \[
      \begin{array}{r@{\,}@{}lr}
        \BRANCH = 
          &
          [
            \neg\tup{@_2\astep[b]\neq_\cmp @_2@_1},
            2{:}\tup{\astep[a]}3,
            \tup{@_3(1\land p)? =_\cmp @_2\astep[b]},
          ] + \dots + \\
          &
          [
            3{:}1,3{:}p,1{:}3,1{:}p, 2{:}\tup{\astep[a]}1,
            \tup{@_3 =_\cmp @_2\astep[b]},
            \tup{@_1 =_\cmp @_2\astep[b]},
            \tup{@_2\astep[b]=_\cmp@_1}
          ]
      \end{array}
    \]
    Notice that $1{:}p$ is obtained from (copy), $2{:}\tup{\child}1$ from (copy$_0$) and $\tup{@_1=_\cmp @_2\astep[b]}$ from (copy$_1$). 
    This way, the nominal $3$ is replaced by the root nominal $1$.
    Finally, the last node expression in $\BRANCH$ is obtained by applying the rule (com$_1$).

  \item[Step 4]
    In the `explore' block of the first recursive call we choose to apply the rule of type 3 with premiss $\tup{@_2\astep[b]=_\cmp@_1}$.
    Before the recursive call we get
    \[
      \Phi =
        [
          \neg\tup{@_2\astep[b]\neq_\cmp @_2@_1}, 1{:}p, 2{:}\tup{\astep[a]}1, 2{:}\tup{\astep[b]}4, \tup{@_4 =_{\cmp} @_1}
        ]
    \]

  \item[Step 5]
    We call $\SATALG(\Phi,\NOM)$ recursively for a second time (with $\NOM=5$).
    In the `expand' block of the second recursive call, we apply all rules of type 1 and get
    \[
      \begin{array}{r@{\,}@{}lr}
        \BRANCH = 
            &
            [
              \neg\tup{@_2\astep[b]\neq_\cmp @_2@_1},
              2{:}\tup{\astep[b]}4,
              \tup{@_4 =_{\cmp} @_2@_1}
            ] + \dots + \\
            &
            [
              \neg \tup{@_4 \neq_{\cmp} @_2@_1}
            ] + \dots +
              & \text{($\lnot$child)} \\
            &
            [
              \neg \tup{@_4 \neq_{\cmp} @_1}
            ]
              & \text{($\neg@$)}
      \end{array}
    \]
        

  \item[Step 6] 
    Since no rules of type 2 are triggered, no clash is found, and there are some untreated node expressions, the algorithm proceeds to the `explore' phase again. A new rule of type 3 is triggered, this time with premise $\psi= \tup{@_2\astep[c]@_1(\neg p)? =_\cmp @_2\astep[b]}$. Before the new recursive call, we get
    \[
      \Phi =
        [
          \neg\tup{@_2\neq_{\cmp}@_1}, 1{:}p
        ] 
        + \dots + 
        [
          2{:}\tup{\astep[c]}5, \tup{@_5@_1(\neg p)? =_\cmp @_2\astep[b]} 
        ]
    \]


  \item[Step 7] 
    The recursive call invokes $\SATALG(\Phi,\NOM)$, with $\NOM=6$. In the `expand' block, we get 
    \[
      \begin{array}{r@{\,}@{}lr}
        \BRANCH = 
            &
            [
              \neg\tup{@_2\neq_{\cmp}@_1}, 1{:}p, 2{:}\tup{\astep[c]}5, \tup{@_5@_1(\neg p)? =_\cmp @_2\astep[b]}
            ] + \dots + \\
            &
            [
              \tup{@_1(\neg p)? =_{\cmp} @_2\astep[b]}, \tup{@_1 =_{\cmp}@_2\astep[b]}, 1{:}\neg p
            ] 
              & (@), (?) 
      \end{array}
    \]
    \item[Step 8] 
      After executing the Step 7, the algorithm finds a clash in line 21, since both $1{:}p$ and $1{:}\neg p$ belong to the sole tableau branch. Therefore, \Cref{alg:sat} returns that $\varphi$ is unsatisfiable.
\end{description}

Notice how in each step the algorithm keeps track of the list of accesibility constraints created in the exploration of diamonds in a DFS manner.
The length of these lists is bounded by the lengths of the path expressions traversed in each diamond.
Moreover, we only need to keep track of a polynomially bounded number of node expressions at each point in these paths.
This example illustrates the polynomial space bound in \Cref{th:satpspace}.
\end{exa}